\documentclass[10pt]{beamer}
\usepackage[utf8]{inputenc}

\usetheme{Madrid}
\usecolortheme{default}
\usepackage{graphicx}
\usepackage{amsmath}
\usepackage{booktabs}
\usepackage{amsmath,amssymb,amsthm,microtype}
\usepackage{tikz}
\graphicspath{{../OverleafPaper/figures/}{OverleafPaper/figures/}{../forgetting/outputs/test/monte_carlo/figures/}{forgetting/outputs/test/monte_carlo/figures/}{../forgetting/outputs/test/first_2000/figures/individual_states/p0_0.5/}{forgetting/outputs/test/first_2000/figures/individual_states/p0_0.5/}{../forgetting/outputs/test/first_2000/figures/state_distributions/p0_0.5/}{forgetting/outputs/test/first_2000/figures/state_distributions/p0_0.5/}}

\title{Weekly report}
\author{Natalia Trigo Tomasevich}
\date{October 11, 2026}

\begin{document}

\frame{\titlepage}

\begin{frame}{This week}
    \begin{itemize}[<+->]
        \item Consolidated results for the forgetting model
        \item Competition: Discussion of modeling choices 
    \end{itemize}
\end{frame}

% Source for all numbered results:
% Progress notes/08_Final_note_intermediate_forgetting/08_final_note_intermediate_forgetting.tex
% The first five results concern the discrete model; the last three concern
% the fluid model and its demand assumptions. Numbering follows the source.

\begin{frame}{Positive advantage: a fixed-policy test}
    Results from \emph{Final note: Intermediate forgetting} on Overleaf.

    \medskip
    Let $U_\rho^i$ be the value of always
    choosing product $i$. On the invariant limit set $\mathcal K_\rho$, define
    \begin{align}
        a_\rho(z)&=-\Delta c+\gamma\Delta p
        [V_\rho(T_\rho^s z)-V_\rho(T_\rho^f z)],\\
        \mathcal A_\rho&=\max_{z\in\mathcal K_\rho}a_\rho(z).
    \end{align}
    
    \begin{block}{Proposition 1: Characterization of positive advantage}
        For $\rho\in(0,1)$,
        \[
        \mathcal A_\rho>0
        \quad\Longleftrightarrow\quad
        \max_{z\in\mathcal K_\rho}
        [U_\rho^1(T_\rho^s z)-U_\rho^1(T_\rho^f z)]
        >\frac{\Delta c}{\gamma\Delta p}.
        \]
    \end{block}
\end{frame}

\begin{frame}{Persistent use of product 2}
    \begin{block}{Proposition 2}
        Fix $\rho\in(0,1)$. If $\mathcal A_\rho>0$, there exists
        $\eta>0$ such that, from every initial state and under any
        optimal policy,
        \[
        \liminf_{N\to\infty}\frac{1}{N}
        \sum_{n=1}^{N}\mathbf 1\{i_n=2\}
        \ge\eta
        \qquad\text{almost surely}.
        \]
        Here $i_n$ is the product chosen at the $n$th interaction
        with provider A.
    \end{block}

    \medskip
    Product 2 has a positive lower asymptotic frequency of use.
    In particular, the provider uses it infinitely often.
\end{frame}

\begin{frame}{State-dependent choices and indifference}
    Define the fixed-policy advantages
    \[
    \begin{aligned}
        a_\rho^i(z)&=-\Delta c+\gamma\Delta p
        [U_\rho^i(T_\rho^s z)-U_\rho^i(T_\rho^f z)],\quad i=1,2,\\
        \mathcal A_\rho^1&=\max_{z\in\mathcal K_\rho}a_\rho^1(z),
        \qquad
        \mathcal A_\rho^2=\min_{z\in\mathcal K_\rho}a_\rho^2(z).
    \end{aligned}
    \]

    \begin{block}{Lemma 1}
        If $\rho\in(1/2,1)$, $\mathcal A_\rho^1>0$, and
        $\mathcal A_\rho^2<0$, each product is strictly optimal at some
        states in $\mathcal K_\rho$. Moreover,
        \[
        \exists\,x^*\in(0,1):\qquad a_\rho(x^*,1-x^*)=0.
        \]
    \end{block}

    \begin{itemize}
        \item Connectedness and continuity guarantee a zero; an interval
        of indifference is also possible.
        \item A threshold structure needs further restrictions on the sign
        pattern: the number, order, and direction of switches remain
        undetermined.
    \end{itemize}
    % Both signs occur, but a zero need not be an isolated crossing:
    % an interval of indifference may separate the strict preference regions.
    % Uniqueness of the zero is sufficient, but is not necessary for a
    % threshold selection. No pointwise equality of fixed-policy and
    % optimal-policy advantages is asserted.
\end{frame}

\begin{frame}{Persistent use of both products}
    \begin{block}{Corollary 1}
        Fix $\rho\in(0,1)$. If $\mathcal A_\rho^1>0$ and
        $\mathcal A_\rho^2<0$, there exist $\eta_1,\eta_2>0$ such that,
        from every initial state and under any optimal policy,
        \[
        \liminf_{N\to\infty}\frac{1}{N}
        \sum_{n=1}^{N}\mathbf 1\{i_n=j\}
        \ge\eta_j
        \quad\text{almost surely},\qquad j=1,2.
        \]
    \end{block}

    \begin{itemize}
        \item Both products are used infinitely often, each with a positive
        lower asymptotic frequency.
    \end{itemize}
\end{frame}

\begin{frame}{\large Cost region with both products strictly optimal at some states}
    \footnotesize
    \setlength{\abovedisplayskip}{4pt}
    \setlength{\belowdisplayskip}{4pt}
    \setlength{\abovedisplayshortskip}{2pt}
    \setlength{\belowdisplayshortskip}{2pt}
    \setlength{\jot}{2pt}
    For $i\in\{1,2\}$, let
    \[
    H_\rho^i(z)
    =\mathbb E_z^i\left[\sum_{t=1}^{\infty}\gamma^{t-1}\mathbf 1\{y_t=A\}\right]
    \]
    be the expected discounted number of purchases from seller~$A$ when
    product~$i$ is always used. Here $t$ counts calendar periods, and
    $\mathbb E_z^i$ denotes expectation starting from $z$ under this fixed policy.

    \begin{block}{Proposition 3}
        Fix $\rho\in(1/2,1)$, $\gamma\in(0,1)$, $p_0\in(0,1)$,
        $0<p_1<p_2<1$, and $0\leq c_1<R$. Define
        \begin{align*}
        \underline c_2
        &=\frac{
         c_1+\gamma\Delta p\,R
         \displaystyle\min_{z\in\mathcal K_\rho}
         \left[H_\rho^2(T_\rho^s z)-H_\rho^2(T_\rho^f z)\right]
        }{
         1+\gamma\Delta p
         \displaystyle\min_{z\in\mathcal K_\rho}
         \left[H_\rho^2(T_\rho^s z)-H_\rho^2(T_\rho^f z)\right]
        },\\[1ex]
        \overline c_2
        &=c_1+\gamma\Delta p\,(R-c_1)
        \max_{z\in\mathcal K_\rho}
        \left[H_\rho^1(T_\rho^s z)-H_\rho^1(T_\rho^f z)\right].
        \end{align*}
        These bounds do not depend on $c_2$. For any $c_2\in(c_1,R)$,
        \begin{equation*}
        \mathcal A_\rho^1>0\ \text{and}\ \mathcal A_\rho^2<0
        \quad\Longleftrightarrow\quad
        \underline c_2<c_2<\overline c_2.
        \end{equation*}
        Thus $(\underline c_2,\overline c_2)\cap(c_1,R)$ is exactly the cost
        region where each product is strictly optimal at some state in
        $\mathcal K_\rho$. For every cost in this region, there exists
        $x^*\in(0,1)$ such that $a_\rho(x^*,1-x^*)=0$.
    \end{block}
\end{frame}

\begin{frame}{Fluid model: the transformed payoff}

    \begin{block}{Proposition 4: Shape of the transformed payoff}
        Assume $D'(x)>0$ and $D'$ is log-concave on $[0,1]$. Then the set $\{x\in[0,1]:F'(x)>0\}$ is an interval, possibly empty or meeting a boundary. If nonempty, $F$ decreases before this interval, increases within it, and decreases after it, with absent segments omitted.
        If it is empty, product 1 is optimal everywhere.
    \end{block}
\end{frame}

\begin{frame}{Fluid model: a single interval of product 2}

    \begin{block}{Theorem 1}
        Suppose, for some $0\le a<m\le1$,
        \[
        F'<0\ \text{on }(0,a)\cup(m,1),
        \qquad F'>0\ \text{on }(a,m),
        \]
        and every interior zero of $F'$ is simple.

        \medskip
        An optimal policy can be selected that uses product 2 on
        at most one interval and product 1 outside it.
        As $x$ increases, the possible nonconstant patterns are
        \[
        1\to2,\qquad 2\to1,\qquad 1\to2\to1.
        \]
    \end{block}

    % At an attractive threshold $x^*\in(p_1,p_2)$, an optimal holding
    % mixture may use product 2 with fraction $(x^*-p_1)/(p_2-p_1)$.

    \medskip
    If $F$ is nonincreasing, product 1 is optimal everywhere;
    if $F$ is nondecreasing, product 2 is optimal everywhere.
\end{frame}

\begin{frame}{Existence of a nonconstant threshold policy}
    \begin{block}{Corollary 2}
        Under the assumptions of Theorem 1, if in addition we have 
        \begin{align*}
        \max_{x\in[0,1]}\{-\Delta c+\lambda\Delta p(U^1_\rho)'(x)\}&>0,\\
        \min_{x\in[0,1]}\{-\Delta c+\lambda\Delta p(U^2_\rho)'(x)\}&<0.
        \end{align*}
        then there is an optimal policy with either one or two interior thresholds, with one of the three patterns in
        Theorem~1.
    \end{block}
\end{frame}

\begin{frame}[plain,s]
    \setbeamertemplate{navigation symbols}{}
    \vfill
    \centering
    {\usebeamerfont{title}\usebeamercolor[fg]{structure}\LARGE
    Competition model\par}
    \vfill
\end{frame}

\begin{frame}{Competition model}
    \begin{itemize}[<+->]
        \item If we consider the model without forgetting with both sellers knowing all the information is not tractable. 
        \item Because it has no forgetting, the state growns infinitely, and since the state now has 4 simensions, the problem becomes impossible to solve unless we truncate the state space. 
        \item The forgetting model gives us the advantage of finite state space.
        \item How much information we want the sellers to know? 
        \item Go to Francisco's proposal 
    \end{itemize}
    
\end{frame}
\end{document}
